%%%%%%%%%%%%%%%%%%%%%%%%%%%%%%%%%%%%%%%%%%%%%%%%%%%%%%%%%%%%%%%%%%%%%%%%%%%
% Chapter: Operational Amplifiers
% Falstad examples: examples/opamps (generated by falstad/circuits/opamps.py)
%%%%%%%%%%%%%%%%%%%%%%%%%%%%%%%%%%%%%%%%%%%%%%%%%%%%%%%%%%%%%%%%%%%%%%%%%%%

% opfig=<s>: scales coordinates and component size by s, but not the text
\tikzset{opfig/.style={scale=#1, /tikz/circuitikz/bipoles/length=#1cm}}

%%%%%%%%%%%%%%%%%%%%%%%%%%%%%%%%%%%%%%%%%%%%%%%%%%%%%%%%%%%%%%%%%%%%%%%%%%%
\section{The Ideal Operational Amplifier}
%%%%%%%%%%%%%%%%%%%%%%%%%%%%%%%%%%%%%%%%%%%%%%%%%%%%%%%%%%%%%%%%%%%%%%%%%%%
\nsl{A bridge circuit (chapter~4) turns a small change of a sensor
resistance into a voltage of a few millivolts. An analog-to-digital
converter needs volts. In between sits the operational amplifier (op-amp):
a differential amplifier with a very high gain that is almost never used on
its own. Together with a few resistors it becomes a buffer, an amplifier
with an exactly defined gain, a summing or a difference amplifier, an
instrumentation amplifier or an active filter (chapter~6). The trick
behind all of these circuits is negative feedback: the resistors, not the
op-amp, define the behaviour.
\newline
In this chapter we first treat the op-amp as an ideal component and derive
the basic circuits with two simple rules. Then we look at the limits of real
op-amps: finite open-loop gain, finite bandwidth, output swing, offset
voltage and bias currents. At the end we design the measurement amplifier
of the lab capstone task ``From Sensor to Digital Value''. Good
introductions are \cite{horowitz2015} (chapter~4), \cite{mancini2002} and
\cite{tietze2008}; \cite{jung2005} collects many practical details.
\newpage}

\begin{description}
\item[{\rot\bf The op-amp: symbol and model}]~\\[-\lblskip]
\begin{center}
\begin{circuitikz}[opfig=1.3]
  \draw (0,0) node[op amp](OA){};
  \draw (OA.-) -- ++(-0.6,0) node[left]{$V_-$ (inverting input)};
  \draw (OA.+) -- ++(-0.6,0) node[left]{$V_+$ (non-inverting input)};
  \draw (OA.out) -- ++(0.6,0) node[right]{$\Vout = A\,(V_+ - V_-)$};
  \draw (OA.up) -- ++(0,0.3) node[vcc]{$+V_S$};
  \draw (OA.down) -- ++(0,-0.3) node[vee]{$-V_S$};
\end{circuitikz}
\end{center}
\bi
	\ite amplifies the \textbf{difference} of its two input voltages
	\ite open-loop gain $A$: typically $10^5 \ldots 10^7$ (\SIrange{100}{140}{\deci\bel})
	\ite supply $\pm V_S$ (or single supply $V_S$ and ground); the output cannot leave the supply range
	\ite supply pins are often omitted in drawings -- but they are always there
\ei
\os{\newpage}

\item[{\rot\bf Properties of the ideal op-amp}]~\\[-\lblskip]
\begin{center}
\begin{tabular}{lll}
\toprule
property & ideal & real (typical) \\
\midrule
open-loop gain $A$ & $\infty$ & $10^5 \ldots 10^7$ \\
input resistance & $\infty$ & \SI{1}{\mega\ohm} \ldots \SI{1}{\tera\ohm} \\
input currents & 0 & \si{\pico\ampere} (CMOS) \ldots \SI{100}{\nano\ampere} (bipolar) \\
output resistance & 0 & \SIrange{10}{100}{\ohm} (open loop) \\
offset voltage & 0 & \SI{10}{\micro\volt} \ldots \SI{5}{\milli\volt} \\
bandwidth & $\infty$ & GBW \SI{1}{\mega\hertz} \ldots \SI{100}{\mega\hertz} \\
\bottomrule
\end{tabular}
\end{center}
\nsl{The ideal values are good approximations in most circuits of this
lecture. Section~\ref{sec:opamp-limits} shows when they are not: high gain,
high frequency, high source resistance and output voltages close to the
supply rails.}
\os{\newpage}

\item[{\rot\bf Open loop: the op-amp as a comparator}]~\\[-\lblskip]
\begin{center}
\begin{tikzpicture}
\begin{axis}[width=12cm, height=5.2cm, xmin=-400, xmax=400, ymin=-16, ymax=16,
  xlabel={$V_+ - V_-$ in \si{\micro\volt}}, ylabel={$\Vout$ in \si{\volt}},
  grid=major, ytick={-13.5,0,13.5}, yticklabel style={font=\scriptsize},
  xticklabel style={font=\scriptsize}, label style={font=\small}]
  \addplot[rwuviolet, very thick] coordinates {(-400,-13.5) (-135,-13.5) (135,13.5) (400,13.5)};
  \addplot[rwucyan, dashed] coordinates {(-400,15) (400,15)};
  \addplot[rwucyan, dashed] coordinates {(-400,-15) (400,-15)};
\end{axis}
\end{tikzpicture}
\end{center}
\bi
	\ite with $A=10^5$ and \SI{\pm 13.5}{\volt} output swing the linear range is only \SI{\pm 135}{\micro\volt}
	\ite any larger difference drives the output to a rail: $\Vout$ tells only \emph{which input is higher}
	\ite this is a \textbf{comparator} -- the building block of the flash ADC (chapter~8)
\ei
\nsl{Without feedback an op-amp is useless as a linear amplifier: a
difference of a few hundred microvolts is enough to saturate it, and the
gain $A$ varies from part to part by a factor of ten and with temperature.
As a comparator it works, although dedicated comparator ICs switch much
faster and have logic-compatible outputs. The ADC chapter treats
comparators in detail.}
\os{\newpage}

\item[{\rot\bf Negative feedback and the golden rules}]~\\[-\lblskip]
\bi
	\ite \textbf{negative feedback}: a part of $\Vout$ is fed back to the \textbf{inverting} input
	\ite if $\Vout$ rises, $V_-$ rises, the difference $V_+-V_-$ shrinks and pulls $\Vout$ back
	\ite the circuit settles where $V_+-V_- = \Vout/A \approx 0$
\ei
\satz{golden}{\textbf{Golden rules} for an op-amp with negative feedback:
\begin{enumerate}
\item The output does whatever is necessary to make the voltage difference
      between the inputs zero: $V_+ = V_-$.
\item The inputs draw no current: $I_+ = I_- = 0$.
\end{enumerate}}
\nsl{The two rules come from \cite{horowitz2015}. Rule 1 holds only as
long as the feedback loop works, i.e.\ the output is not saturated at a
rail and the feedback goes to the inverting input. With feedback to the
non-inverting input (positive feedback) the circuit becomes a Schmitt
trigger or an oscillator. Rule 2 is a property of the op-amp itself and
holds always (within the bias current).}
\os{\newpage}
\end{description}

%%%%%%%%%%%%%%%%%%%%%%%%%%%%%%%%%%%%%%%%%%%%%%%%%%%%%%%%%%%%%%%%%%%%%%%%%%%
\section{Basic Amplifier Circuits}
%%%%%%%%%%%%%%%%%%%%%%%%%%%%%%%%%%%%%%%%%%%%%%%%%%%%%%%%%%%%%%%%%%%%%%%%%%%
\nsl{With the golden rules the analysis of the basic circuits takes only a
few lines. All circuits in this section exist as Falstad examples. Falstad
uses an op-amp with $A=10^5$ and an output clamped to the supply, which is
close enough to ideal for these circuits.\newpage}

\begin{description}
\item[{\rot\bf Voltage follower (buffer)}]~\\[-\lblskip]
\begin{center}
\begin{circuitikz}[opfig=1.3]
  \draw (0,0) node[op amp, noinv input up](OA){};
  \draw (OA.+) to[short,-o] ++(-1,0) node[left]{$\Vin$};
  \draw (OA.-) -- ++(0,-0.9) -| ($(OA.out)+(0.4,0)$);
  \draw (OA.out) to[short,-o] ++(1.2,0) node[right]{$\Vout = \Vin$};
\end{circuitikz}
\end{center}
\bi
	\ite rule 1: $V_- = V_+$, and $V_- = \Vout$ $\Rightarrow$ $\Vout = \Vin$, gain 1
	\ite input current $\approx 0$: the source is \textbf{not loaded}
	\ite output: low resistance, can drive \si{\milli\ampere} into a load
\ei
\nsl{The follower does not amplify the voltage, it amplifies the current: it
copies a voltage from a high-impedance node to a low-impedance output.
This is exactly what a bridge needs. A bridge tap is a voltage divider with
a source resistance of $R_1\parallelto R_2$ (chapter~1); a following stage
with a finite input resistance would load it and shift the measured
voltage. The same holds for the sample-and-hold circuit (chapter~7): the
voltage on the hold capacitor must not be discharged by the following
stage, so it is read by a follower with a very small input current.}
\os{\newpage}

\item[{\rot\bf Why buffering matters}]~\\[-\lblskip]
\bi
	\ite divider \SI{100}{\kilo\ohm}/\SI{100}{\kilo\ohm} at \SI{10}{\volt}, load $R_L=\SI{10}{\kilo\ohm}$
	\ite directly: $\Vout = \SI{10}{\volt}\cdot\frac{R_2\parallelto R_L}{R_1+R_2\parallelto R_L} = \SI{0.833}{\volt}$ instead of \SI{5}{\volt}
	\ite behind a follower: \SI{5.00}{\volt}, the divider does not notice the load
	\ite uses in this course
	\itee bridge taps \(\rightarrow\) instrumentation amplifier inputs
	\itee reference ladder of the flash ADC, hold capacitor of the S\&H
\ei
\falstad{opamp_follower}\par
\os{\newpage}

\item[{\rot\bf Non-inverting amplifier}]~\\[-\lblskip]
\begin{center}
\begin{circuitikz}[opfig=1.3]
  \draw (0,0) node[op amp, noinv input up](OA){};
  \coordinate (X) at ($(OA.out)+(0.5,0)$);
  \draw (OA.+) to[short,-o] ++(-1,0) node[left]{$\Vin$};
  \draw (OA.-) -- ++(0,-1) coordinate(N);
  \draw (N) to[R, l_=$R_2$] (N -| X) -- (X);
  \draw (N) to[R, l_=$R_1$] ++(0,-1.6) node[ground]{};
  \draw (OA.out) -- (X) to[short,-o] ++(0.8,0) node[right]{$\Vout$};
  \draw (N) node[circ]{} node[left]{$V_-$};
\end{circuitikz}
\end{center}
\bi
	\ite $R_1, R_2$ form a divider: $V_- = \Vout\,\frac{R_1}{R_1+R_2}$; rule 1: $V_-=\Vin$
\ei
\keyeq{Non-inverting amplifier}{$\displaystyle G = \frac{\Vout}{\Vin} = 1 + \frac{R_2}{R_1}$
\quad input resistance $\approx\infty$}
\os{\newpage}

\item[{\rot\bf Inverting amplifier}]~\\[-\lblskip]
\begin{center}
\begin{circuitikz}[opfig=1.3]
  \draw (0,0) node[op amp](OA){};
  \coordinate (X) at ($(OA.out)+(0.5,0)$);
  \draw (OA.+) -- ++(0,-0.4) node[ground]{};
  \draw (OA.-) -- ++(-0.4,0) coordinate(N);
  \draw (N) to[R, l_=$R_1$, -o] ++(-2.2,0) node[left]{$\Vin$};
  \draw (N) -- ++(0,1.1) coordinate(T) to[R, l=$R_2$] (T -| X) -- (X);
  \draw (OA.out) -- (X) to[short,-o] ++(0.8,0) node[right]{$\Vout$};
  \draw (N) node[circ]{};
\end{circuitikz}
\end{center}
\bi
	\ite rule 1: $V_- = V_+ = 0$, the inverting input is a \textbf{virtual ground}
	\ite rule 2: the current $\Vin/R_1$ continues through $R_2$: $\Vout = -R_2\cdot \Vin/R_1$
	\ite input resistance $= R_1$ (it loads the source!)
\ei
\keyeq{Inverting amplifier}{$\displaystyle G = \frac{\Vout}{\Vin} = -\frac{R_2}{R_1}$}
\os{\newpage}

\item[{\rot\bf Non-inverting and inverting in Falstad}]~\\[-\lblskip]
\begin{center}
\begin{tikzpicture}
\begin{axis}[width=12.5cm, height=5.2cm, xmin=0, xmax=2, ymin=-1.8, ymax=1.8,
  xlabel={$t$ in \si{\milli\second}}, ylabel={$V$ in \si{\volt}}, grid=major,
  legend pos=outer north east, legend style={font=\scriptsize},
  ticklabel style={font=\scriptsize}, label style={font=\small}]
  \addplot[rwugray, thick, domain=0:2, samples=100] {0.5*sin(360*x)};
  \addplot[rwuviolet, very thick, domain=0:2, samples=100] {1.5*sin(360*x)};
  \addplot[rwucyan, very thick, domain=0:2, samples=100] {-1.0*sin(360*x)};
  \legend{$\Vin$, $G=3$, $G=-2$}
\end{axis}
\end{tikzpicture}
\end{center}
\bi
	\ite $\Vin$: \SI{0.5}{\volt} amplitude, \SI{1}{\kilo\hertz}; $R_1=\SI{10}{\kilo\ohm}$, $R_2=\SI{20}{\kilo\ohm}$ in both circuits
	\ite non-inverting: $G=1+2=3$, \SI{1.5}{\volt} in phase; inverting: $G=-2$, \SI{1.0}{\volt}, phase \ang{180}
\ei
\falstad{opamp_noninv_inv}\par
\os{\newpage}

\item[{\rot\bf Summing amplifier}]~\\[-\lblskip]
\begin{center}
\begin{circuitikz}[opfig=1.3]
  \draw (0,0) node[op amp](OA){};
  \coordinate (X) at ($(OA.out)+(0.5,0)$);
  \draw (OA.+) -- ++(0,-0.4) node[ground]{};
  \draw (OA.-) -- ++(-0.4,0) coordinate(N);
  \draw (N) -- ++(0,1.8) coordinate(T) to[R, l=$R_f$] (T -| X) -- (X);
  \draw (N) -- ++(0,1.0) coordinate(N1) to[R, l_=$R_1$, -o] ++(-2.4,0) node[left]{$V_1$};
  \draw (N) to[R, l_=$R_2$, -o] ++(-2.4,0) node[left]{$V_2$};
  \draw (N) -- ++(0,-1.0) coordinate(N3) to[R, l_=$R_3$, -o] ++(-2.4,0) node[left]{$V_3$};
  \draw (OA.out) -- (X) to[short,-o] ++(0.8,0) node[right]{$\Vout$};
  \draw (N) node[circ]{} (N1) node[circ]{} (N3) node[circ]{};
\end{circuitikz}
\end{center}
\bi
	\ite the virtual ground decouples the inputs; the currents add up in $R_f$
\ei
\keyeq{Inverting summing amplifier}{$\displaystyle \Vout = -R_f\left(\frac{V_1}{R_1}+\frac{V_2}{R_2}+\frac{V_3}{R_3}\right)$}
\nsl{Because the summing node is held at 0\,V, every input sees only its own
resistor and the inputs do not influence each other. Weighted sums are
easy: the weight of input $k$ is $-R_f/R_k$. Applications are audio
mixers, binary weighted DACs and adding an offset voltage to a signal
(level shifting). The Falstad example adds $V_1=\SI{1}{\volt}$ through
\SI{10}{\kilo\ohm}, $V_2=\SI{2}{\volt}$ through \SI{20}{\kilo\ohm} and
$V_3=\SI{-0.5}{\volt}$ through \SI{10}{\kilo\ohm} with $R_f=\SI{10}{\kilo\ohm}$
(exercise~\ref{ex:opamp-summing}).}
\os{\newpage}
\end{description}

%%%%%%%%%%%%%%%%%%%%%%%%%%%%%%%%%%%%%%%%%%%%%%%%%%%%%%%%%%%%%%%%%%%%%%%%%%%
\section{Difference and Instrumentation Amplifiers}
%%%%%%%%%%%%%%%%%%%%%%%%%%%%%%%%%%%%%%%%%%%%%%%%%%%%%%%%%%%%%%%%%%%%%%%%%%%
\nsl{A bridge delivers its signal as the difference of two voltages that both
sit at about half the bridge supply. We need an amplifier that amplifies
the small difference (the differential-mode signal) and ignores the large
common part (the common-mode signal). The difference amplifier does this
with one op-amp and four resistors; the instrumentation amplifier adds
buffers so that the bridge is not loaded.\newpage}

\begin{description}
\item[{\rot\bf Differential and common-mode signal}]~\\[-\lblskip]
\bi
	\ite two input voltages $V_1$, $V_2$ can be written as
	\itee differential-mode voltage $V_D = V_1 - V_2$ \quad (the measurement signal)
	\itee common-mode voltage $V_{CM} = (V_1 + V_2)/2$ \quad (not wanted)
	\ite symmetric bridge at \SI{5}{\volt}: $V_D$ some \si{\milli\volt}, $V_{CM}\approx\SI{2.5}{\volt}$; capstone bridge: $V_D = \SIrange{0}{37}{\milli\volt}$, $V_{CM}\approx\SI{0.1}{\volt}$
	\ite ideal: $\Vout = A_D\,V_D$; real: $\Vout = A_D\,V_D + A_{CM}\,V_{CM}$
\ei
\defn{cmrr}{The \textbf{common-mode rejection ratio} is the ratio of
differential gain to common-mode gain:
$\mathrm{CMRR} = A_D/A_{CM}$, in decibels $20\log_{10}(A_D/A_{CM})$.}
\os{\newpage}

\item[{\rot\bf Difference amplifier}]~\\[-\lblskip]
\begin{center}
\begin{circuitikz}[opfig=1.3]
  \draw (0,0) node[op amp](OA){};
  \coordinate (X) at ($(OA.out)+(0.5,0)$);
  \draw (OA.-) -- ++(-0.4,0) coordinate(N);
  \draw (N) to[R, l_=$R_1$, -o] ++(-2.2,0) node[left]{$V_2$};
  \draw (N) -- ++(0,1.1) coordinate(T) to[R, l=$R_2$] (T -| X) -- (X);
  \draw (OA.+) -- ++(-0.4,0) coordinate(P);
  \draw (P) to[R, l=$R_3$, -o] ++(-2.2,0) node[left]{$V_1$};
  \draw (P) to[R, l=$R_4$] ++(0,-1.6) node[ground]{};
  \draw (OA.out) -- (X) to[short,-o] ++(0.8,0) node[right]{$\Vout$};
  \draw (N) node[circ]{} (P) node[circ]{};
\end{circuitikz}
\end{center}
\bi
	\ite superposition: $\displaystyle\Vout = V_1\,\frac{R_4}{R_3+R_4}\left(1+\frac{R_2}{R_1}\right) - V_2\,\frac{R_2}{R_1}$
	\ite with $R_4/R_3 = R_2/R_1$: $\displaystyle\Vout = \frac{R_2}{R_1}\,(V_1-V_2)$
\ei
\nsl{The derivation uses superposition. With $V_1 = 0$ the circuit is an
inverting amplifier for $V_2$ (the $+$ input sits at 0\,V through $R_3$ and
$R_4$): $\Vout = -V_2 R_2/R_1$. With $V_2=0$ the voltage $V_1$ is divided by
$R_3,R_4$ and then amplified by a non-inverting amplifier with gain
$1+R_2/R_1$. If the two ratios are equal, the factor in front of $V_1$
becomes $\frac{R_2/R_1}{1+R_2/R_1}\,(1+R_2/R_1) = R_2/R_1$ and only the
difference remains. The input resistances are not equal ($R_1$ for $V_2$,
$R_3+R_4$ for $V_1$) and not high: the circuit loads a bridge.}
\os{\newpage}

\item[{\rot\bf Resistor matching limits the CMRR}]~\\[-\lblskip]
\bi
	\ite the common-mode rejection depends only on how well $R_4/R_3$ equals $R_2/R_1$
	\ite example $G=10$, $V_1=\SI{2.505}{\volt}$, $V_2=\SI{2.500}{\volt}$: ideal $\Vout=\SI{50.0}{\milli\volt}$
	\ite $R_4$ only \SI{1}{\percent} too high: $\Vout=\SI{72.6}{\milli\volt}$ -- an error of \SI{45}{\percent}!
	\ite worst case with resistor tolerance $\varepsilon$: $\mathrm{CMRR} \approx \dfrac{1+G}{4\varepsilon}$
	\itee \SI{1}{\percent} resistors, $G=10$: CMRR $\approx 275$ (\SI{49}{\deci\bel})
	\itee remedy: \SI{0.1}{\percent} resistors, matched resistor networks, a trimmer, or an integrated difference amplifier (e.g.\ INA105, AD8276)
\ei
\falstad{opamp_diffamp}\par
\nsl{The common-mode voltage of a symmetric bridge (2.5\,V) is more than a hundred times
larger than its signal (some 10\,mV), so a small common-mode gain already gives
a large error. The worst case formula follows if all four resistors deviate
by $\pm\varepsilon$ in the unfavourable direction \cite{jung2005}. An
op-amp has its own CMRR (typically 80 to 120\,dB), but in a discrete
difference amplifier the resistors dominate. A higher gain $G$ improves the
CMRR, because the differential gain grows while the common-mode error
stays the same.}
\os{\newpage}

\item[{\rot\bf Instrumentation amplifier (three op-amps)}]~\\[-\lblskip]
\begin{center}
\begin{circuitikz}[opfig=1.0]
  % input stage
  \draw (0,2.4) node[op amp, noinv input up](A1){};
  \draw (0,-2.4) node[op amp](A2){};
  \coordinate (yA) at (0,0.8);  \coordinate (yB) at (0,-0.8);
  \coordinate (X1) at ($(A1.out)+(0.5,0)$);
  \coordinate (X2) at ($(A2.out)+(0.5,0)$);
  \draw (A1.+) to[short,-o] ++(-0.8,0) node[left]{$V_2$};
  \draw (A2.+) to[short,-o] ++(-0.8,0) node[left]{$V_1$};
  \draw (A1.-) -- (A1.- |- yA) coordinate(a);
  \draw (A2.-) -- (A2.- |- yB) coordinate(b);
  \draw (a) to[R, l_=$R_G$] (b);
  \draw (a) -- (a -| X1) to[R, l_=$R$] (X1);
  \draw (b) -- (b -| X2) to[R, l=$R$] (X2);
  \draw (A1.out) -- (X1) -- ++(0.8,0) coordinate(o1);
  \draw (A2.out) -- (X2) -- ++(0.8,0) coordinate(o2);
  \draw (a) node[circ]{} (b) node[circ]{} (X1) node[circ]{} (X2) node[circ]{};
  % difference stage
  \draw (7,0) node[op amp](A3){};
  \coordinate (X) at ($(A3.out)+(0.5,0)$);
  \draw (A3.-) -- ++(-0.4,0) coordinate(n);
  \draw (A3.+) -- ++(-0.4,0) coordinate(p);
  \draw (o1) |- ($(n)+(-2,0)$) to[R, l=$R_1$] (n);
  \draw (o2) |- ($(p)+(-2,0)$) to[R, l_=$R_3$] (p);
  \draw (p) to[R, l=$R_4$] ++(0,-1.8) node[ground]{};
  \draw (n) -- ++(0,1.2) coordinate(t) to[R, l=$R_2$] (t -| X) -- (X);
  \draw (A3.out) -- (X) to[short,-o] ++(0.8,0) node[right]{$\Vout$};
  \draw (n) node[circ]{} (p) node[circ]{};
  \node[above left, font=\scriptsize] at (A1.out) {A1}; \node[below left, font=\scriptsize] at (A2.out) {A2}; \node[above left, font=\scriptsize] at (A3.out) {A3};
\end{circuitikz}
\end{center}
\keyeq{Instrumentation amplifier}{$\displaystyle \Vout = \left(1+\frac{2R}{R_G}\right)\frac{R_2}{R_1}\,(V_1 - V_2)$}
\os{\newpage}

\item[{\rot\bf How the three op-amp INA works}]~\\[-\lblskip]
\bi
	\ite A1, A2 are non-inverting: the inputs draw no current, the bridge is not loaded
	\ite rule 1 puts $V_1$ and $V_2$ across $R_G$: current $(V_1-V_2)/R_G$ flows through $R, R_G, R$
	\ite difference at the outputs of A1/A2: $(V_1-V_2)(1+2R/R_G)$; the common mode passes with gain 1
	\ite A3 is a difference amplifier with gain $R_2/R_1$ and removes the common mode
	\ite one resistor $R_G$ sets the gain; integrated INAs (INA128, AD620, AD8221) contain everything except $R_G$
\ei
\nsl{The input stage amplifies the differential signal by $1+2R/R_G$ but
the common-mode signal only by one. This improves the CMRR of the
complete amplifier by the factor $1+2R/R_G$ compared to the difference
stage alone. Integrated instrumentation amplifiers have laser-trimmed
resistors and reach a CMRR of 100\,dB and more. Their reference pin (the
lower end of $R_4$) shifts the output voltage, which is used for level
shifting and offset trimming (section~\ref{sec:opamp-design}).}
\os{\newpage}

\item[{\rot\bf The lab version: two buffers and a difference amplifier}]~\\[-\lblskip]
\begin{center}
\begin{circuitikz}[opfig=0.95]
  % bridge
  \draw (0,0) to[R, l=$R_S$] (0,2) to[R, l=$R_1$] (0,4) -- (2,4) to[R, l=$R_3$] (2,2) to[R, l=$R_4$] (2,0) -- (0,0);
  \draw (1,4) node[vcc]{\SI{4.00}{\volt}};
  \draw (1,0) node[ground]{};
  \draw (0,2) to[short, *-o] ++(-0.8,0) node[left]{$V_1$};
  \draw (2,2) to[short, *-o] ++(0.8,0) node[right]{$V_2$};
  % buffers
  \draw (6,3) node[op amp, noinv input up](B1){};
  \draw (6,-1) node[op amp, noinv input up](B2){};
  \draw (B1.+) to[short,-o] ++(-0.8,0) node[left]{$V_2$};
  \draw (B2.+) to[short,-o] ++(-0.8,0) node[left]{$V_1$};
  \draw (B1.-) -- ++(0,-0.6) -| ($(B1.out)+(0.3,0)$);
  \draw (B2.-) -- ++(0,-0.6) -| ($(B2.out)+(0.3,0)$);
  % difference amplifier
  \draw (11,1) node[op amp](D){};
  \coordinate (X) at ($(D.out)+(0.5,0)$);
  \draw (D.-) -- ++(-0.3,0) coordinate(n);
  \draw (D.+) -- ++(-0.3,0) coordinate(p);
  \draw (B1.out) -- ++(0.6,0) |- ($(n)+(-2,0)$) to[R, l=$R_1$] (n);
  \draw (B2.out) -- ++(0.6,0) |- ($(p)+(-2,0)$) to[R, l_=$R_3$] (p);
  \draw (p) to[R, l=$R_4$] ++(0,-1.8) node[ground]{};
  \draw (n) -- ++(0,1.1) coordinate(t) to[R, l=$R_2$] (t -| X) -- (X);
  \draw (D.out) -- (X) to[short,-o] ++(0.8,0) node[right]{$V_\mathrm{diff}$};
  \draw (n) node[circ]{} (p) node[circ]{};
\end{circuitikz}
\end{center}
\bi
	\ite followers isolate the bridge; the difference stage sets the gain $G_1=R_2/R_1$
	\ite simpler than the three op-amp version, but all gain must come from $R_2/R_1$
\ei
\nsl{This is the circuit of the lab: the two buffers make the input
resistance practically infinite, so the bridge voltages are not changed by
the difference amplifier. Compared to the three op-amp INA the input stage
has gain one, so the CMRR is determined by the matching of $R_1 \ldots R_4$
alone. In the capstone circuit the gain is distributed over the difference
stage and a following non-inverting stage (section~\ref{sec:opamp-design}).}
\os{\newpage}
\end{description}

%%%%%%%%%%%%%%%%%%%%%%%%%%%%%%%%%%%%%%%%%%%%%%%%%%%%%%%%%%%%%%%%%%%%%%%%%%%
\section{Limits of Real Op-Amps}\label{sec:opamp-limits}
%%%%%%%%%%%%%%%%%%%%%%%%%%%%%%%%%%%%%%%%%%%%%%%%%%%%%%%%%%%%%%%%%%%%%%%%%%%
\nsl{The golden rules assume an infinite open-loop gain, zero input
currents and an output that can reach any voltage. A measurement amplifier
for a bridge works with a total gain of about 100 and an input signal of a
few tens of millivolts. Here the deviations of a real op-amp become
visible: the gain is a little lower than calculated, the output shows a
voltage although the input is zero, and the output cannot reach the supply
rails. This section quantifies these effects and shows how to design
around them \cite{mancini2002, horowitz2015}.\newpage}

\begin{description}
\item[{\rot\bf Finite open-loop gain}]~\\[-\lblskip]
\begin{center}
\begin{tikzpicture}[font=\small, >=Stealth, box/.style={draw, thick, minimum width=12mm, minimum height=8mm}]
  \node[circle, draw, inner sep=1.5pt] (s) at (0,0) {$\Sigma$};
  \node[box, right=12mm of s] (a) {$A$};
  \node[box, below=6mm of a] (b) {$\beta = \frac{1}{G}$};
  \draw[->] (-1.6,0) node[left]{$\Vin$} -- (s);
  \node[above left] at (s.west) {$+$};
  \draw[->] (s) -- node[above, font=\scriptsize]{$V_+-V_-$} (a);
  \draw[->] (a) -- (4.2,0) node[right]{$\Vout$};
  \draw[->] (3.4,0) |- (b);
  \draw[->] (b) -| (s);
  \node[below left] at (s.south) {$-$};
\end{tikzpicture}
\end{center}
\bi
	\ite $\Vout = A\,(\Vin - \beta\,\Vout)$, with the feedback factor $\beta = 1/G$ (non-inverting: $\beta = R_1/(R_1+R_2)$)
	\ite loop gain $A\beta = A/G$: the larger, the closer the circuit comes to the ideal gain $G$
\ei
\keyeq{Closed-loop gain with finite open-loop gain $A$}{$\displaystyle G_\mathrm{real} = \frac{A}{1+A\beta} = \frac{G}{1+G/A} = \frac{1}{1/G + 1/A}$
\qquad relative error $\displaystyle\approx -\frac{G}{A}$}
\os{\newpage}

\item[{\rot\bf How large is the error?}]~\\[-\lblskip]
\begin{center}
\begin{tabular}{rrrr}
\toprule
ideal gain $G$ & $G_\mathrm{real}$ ($A=10^5$) & error & $\Vout$ for \SI{1}{\milli\volt} \\
\midrule
10 & 9.999 & \SI{-0.01}{\percent} & \SI{10.0}{\milli\volt} \\
100 & 99.90 & \SI{-0.1}{\percent} & \SI{99.9}{\milli\volt} \\
1000 & 990.1 & \SI{-1}{\percent} & \SI{0.990}{\volt} \\
10000 & 9091 & \SI{-9.1}{\percent} & \SI{9.09}{\volt} \\
\bottomrule
\end{tabular}
\end{center}
\bi
	\ite the error grows proportionally to $G$: one stage with $G>1000$ shows a visible error
	\ite \textbf{split the gain}: $G = 10\cdot 100$ gives 998.9 (\SI{-0.11}{\percent}), $G=100\cdot 100$ gives 9980 (\SI{-0.2}{\percent})
	\ite rule of the lab: aim for $G<100$ per stage
\ei
\falstad{opamp_finite_gain}\par
\nsl{The relative errors of cascaded stages add up, but each stage has a
small error because its own gain is small: $10/A + 100/A = 110/A$ instead
of $1000/A$. Note that $A$ itself is not a precise number: the data sheet
gives a minimum, and $A$ decreases with temperature, with the load and,
most importantly, with frequency. A circuit that relies on a loop gain of
only a few is therefore not precise, even if the formula above corrects
for it.}
\os{\newpage}

\item[{\rot\bf Gain-bandwidth product}]~\\[-\lblskip]
\begin{center}
\begin{tikzpicture}
\begin{semilogxaxis}[width=12.5cm, height=5.4cm, xmin=1, xmax=1e7, ymin=-10, ymax=110,
  xlabel={$f$ in \si{\hertz}}, ylabel={gain in \si{\deci\bel}}, grid=major,
  ytick={0,20,40,60,80,100}, ticklabel style={font=\scriptsize}, label style={font=\small},
  legend style={font=\scriptsize, at={(0.98,0.95)}, anchor=north east}]
  \addplot[rwuviolet, very thick, domain=0:7, samples=120] ({10^x}, {20*log10(1e5/sqrt(1+(10^x/10)^2))});
  \addplot[rwucyan, very thick, domain=0:7, samples=120] ({10^x}, {20*log10(100/sqrt(1+(10^x/1e4)^2))});
  \addplot[rwugray, dashed] coordinates {(1e4,-10) (1e4,40)};
  \legend{open loop $A(f)$, closed loop $G=100$}
\end{semilogxaxis}
\end{tikzpicture}
\end{center}
\bi
	\ite above a low corner frequency (here \SI{10}{\hertz}) $A$ falls by \SI{20}{\deci\bel} per decade: $A(f)\approx \mathrm{GBW}/f$
	\ite closed-loop bandwidth $f_c \approx \mathrm{GBW}/G$ (here \SI{1}{\mega\hertz}/100 = \SI{10}{\kilo\hertz})
\ei
\os{\newpage}

\item[{\rot\bf Consequences of the GBW}]~\\[-\lblskip]
\bi
	\ite on the falling slope gain and frequency are inversely proportional:
\ei
\keyeq{Gain-bandwidth product}{$\displaystyle \mathrm{GBW} = A(f)\cdot f = G\cdot f_c$ \qquad (constant above the corner frequency)}
\bi
	\ite at \SI{1}{\kilo\hertz} an op-amp with GBW = \SI{1}{\mega\hertz} has only $A=1000$: a $G=100$ stage then has \SI{10}{\percent} error
	\ite splitting the gain also helps here: two stages $G=10$ each have \SI{100}{\kilo\hertz} bandwidth
	\ite our sensor signal is slow (a few \si{\hertz}): the GBW is no problem in the capstone
	\ite \textbf{Falstad does not simulate the GBW}: its op-amp has a constant $A=10^5$ at all frequencies
\ei
\nsl{The open-loop gain of almost every op-amp is deliberately rolled off
by an internal capacitor (dominant-pole compensation), so that the
amplifier stays stable with any amount of negative feedback. The result is
a first-order low-pass with a very low corner frequency. The product of
gain and frequency is constant on the falling slope and is listed in the
data sheet as gain-bandwidth product (GBW) or unity-gain frequency $f_T$.
For the closed loop the gain that counts is the noise gain $1/\beta$, which
is $1+R_2/R_1$ for both the non-inverting and the inverting amplifier.
A related limit for large signals is the slew rate, the maximum rate of
change of the output voltage (e.g.\ \SI{0.5}{\volt/\micro\second}).
In Falstad and in our simulator the op-amp has no bandwidth limit, so
results of fast circuits must be checked against the data sheet.}
\os{\newpage}

\item[{\rot\bf Output swing and rail-to-rail}]~\\[-\lblskip]
\bi
	\ite the output cannot exceed the supply; a classic op-amp stays \SIrange{1}{2}{\volt} away from the rails
	\itee LM358 at \SI{5}{\volt}: output from about \SI{0}{\volt} to \SI{3.5}{\volt}
	\ite \textbf{rail-to-rail output}: reaches the rails within some \SI{10}{\milli\volt} (light load) to a few \SI{100}{\milli\volt} (\si{\milli\ampere} load)
	\ite the \textbf{input common-mode range} is limited as well: not every op-amp works with inputs near the rails (\textbf{rail-to-rail input}); the capstone bridge sits at only \SI{0.1}{\volt}
	\ite consequence for the capstone: an output of exactly \SI{0}{\volt} is not possible on a single \SI{5}{\volt} supply -- use a negative supply, or a small offset and a range of e.g.\ \SIrange{0.1}{4.1}{\volt}
\ei
\nsl{When the output hits a rail the feedback loop is broken, rule~1 no
longer holds, and the signal is clipped. Clipping cannot be repaired by
later stages: a clipped 50\,Hz interference becomes a DC error after the
low-pass filter. Each stage therefore needs headroom for the signal plus
all interference it has to carry. Falstad's op-amp model clamps exactly at
the values \texttt{vmax} and \texttt{vmin}, which behaves like an ideal
rail-to-rail output.}
\os{\newpage}

\item[{\rot\bf Input offset voltage and bias currents}]~\\[-\lblskip]
\begin{center}
\begin{circuitikz}[opfig=1.2]
  \draw (0,0) node[op amp, noinv input up](OA){};
  \coordinate (X) at ($(OA.out)+(0.5,0)$);
  \draw (OA.+) -- ++(-0.3,0) coordinate(p) to[V, invert, l_=$V_{OS}$] ++(-1.6,0) coordinate(q) to[R, l_=$R_S$] ++(-1.8,0) -- ++(0,-0.6) node[ground]{};
  \draw (q) to[I, l_=$I_B$] ++(0,-1.4) node[ground]{};
  \draw (q) node[circ]{};
  \draw (OA.-) -- ++(0,-0.8) coordinate(N);
  \draw (N) to[R, l_=$R_2$] (N -| X) -- (X);
  \draw (N) to[R, l_=$R_1$] ++(0,-1.4) node[ground]{};
  \draw (OA.out) -- (X) to[short,-o] ++(0.8,0) node[right]{$\Vout$};
\end{circuitikz}
\end{center}
\bi
	\ite offset $V_{OS}$: a small voltage source in series with an input, amplified by $1+R_2/R_1$
	\ite bias current $I_B$ flows through the source resistance $R_S$ and adds $R_S I_B$
	\ite $\Vout = (1 + R_2/R_1)\,(V_{OS} + R_S I_B)$ with $\Vin = 0$
\ei
\os{\newpage}

\item[{\rot\bf Offset after a high gain}]~\\[-\lblskip]
\bi
	\ite typical values
	\itee general purpose bipolar (LM358): $V_{OS}\approx\SI{2}{\milli\volt}$, $I_B\approx\SI{20}{\nano\ampere}$
	\itee CMOS (MCP6002, TLV9002): $V_{OS}\approx\SI{1}{\milli\volt}$, $I_B\approx\SI{1}{\pico\ampere}$
	\itee zero-drift / chopper (OPA2333, MCP6V02): $V_{OS}<\SI{10}{\micro\volt}$
	\ite Falstad example with $G=101$: $V_{OS}=\SI{1}{\milli\volt}$ gives \SI{0.101}{\volt}; $I_B=\SI{100}{\nano\ampere}$ through \SI{100}{\kilo\ohm} gives \SI{-1.01}{\volt}
	\ite capstone: \SI{2}{\milli\volt} offset $\times$ 107.6 = \SI{0.22}{\volt} -- \SI{5}{\percent} of the full scale, almost half an LSB!
	\ite remedies: zero-drift op-amps, CMOS inputs, low source resistances, \textbf{offset trimming}
\ei
\falstad{opamp_offset}\par
\nsl{The offset voltage is a property of the input transistor pair: the
two transistors are never perfectly equal. It is specified as a typical
and a maximum value and it drifts with temperature (typically
1 to 10\,\textmu V/K). Because it appears at the input, it is amplified by
the full noise gain. A bridge signal of 37\,mV and an offset of 2\,mV make
an error of more than 5\,\%. Bias currents are the base currents (bipolar)
or gate leakage currents (CMOS) of the input transistors. In a bridge
amplifier the source resistance is low (about 100\,$\Omega$ for the
capstone bridge), so $I_B$ matters little; but behind a sample-and-hold
capacitor or a high-ohmic filter it does. The Falstad model of the op-amp
has no offset and no bias current; the example models them with an extra
voltage source and current source.}
\os{\newpage}

\item[{\rot\bf Single-supply operation and level shifting}]~\\[-\lblskip]
\begin{center}
\begin{circuitikz}[opfig=1.15]
  \draw (0,0) node[op amp, noinv input up](OA){};
  \coordinate (X) at ($(OA.out)+(0.5,0)$);
  \draw (OA.up) -- ++(0,0.3) node[vcc]{\SI{5}{\volt}};
  \draw (OA.down) -- ++(0,-0.3) node[ground]{};
  \draw (OA.+) -- ++(-0.4,0) coordinate(p);
  \draw (p) -- ++(0,1.0) coordinate(p1) to[R, l_=$R$, -o] ++(-2,0) node[left]{$\Vin$};
  \draw (p) to[R, l=$R$, -o] ++(-2,0) node[left]{$\Vref = \SI{2.5}{\volt}$};
  \draw (OA.-) -- ++(0,-1.3) coordinate(N);
  \draw (N) to[R, l_=$R$] (N -| X) -- (X);
  \draw (N) to[R, l_=$R$] ++(0,-1.3) node[ground]{};
  \draw (OA.out) -- (X) to[short,-o] ++(0.8,0) node[right]{$\Vout = \Vin + \Vref$};
  \draw (p) node[circ]{};
\end{circuitikz}
\end{center}
\bi
	\ite single supply: the output can only be positive; bipolar signals need a \textbf{bias point} (often $V_S/2$)
	\ite here $V_+ = (\Vin+\Vref)/2$, gain 2: $\Vin=\pm\SI{2}{\volt}$ becomes \SIrange{0.5}{4.5}{\volt}
	\ite a follower on the same supply clips the negative half wave at \SI{0}{\volt}
\ei
\falstad{opamp_level_shift}\par
\os{\newpage}
\end{description}

%%%%%%%%%%%%%%%%%%%%%%%%%%%%%%%%%%%%%%%%%%%%%%%%%%%%%%%%%%%%%%%%%%%%%%%%%%%
\section{Designing the Measurement Amplifier}\label{sec:opamp-design}
%%%%%%%%%%%%%%%%%%%%%%%%%%%%%%%%%%%%%%%%%%%%%%%%%%%%%%%%%%%%%%%%%%%%%%%%%%%
\nsl{Now we design the front end of the capstone task (chapter~9). The
sensor is a Pt100 (\SI{100}{\ohm} at \SI{0}{\celsius}, \SI{138.5}{\ohm} at
\SI{100}{\celsius}) in a Wheatstone bridge (chapter~4). The bridge is fed
from the same \SI{4.00}{\volt} reference that later supplies the ADC ladder,
with $R_1 = R_3 = \SI{3.9}{\kilo\ohm}$ and $R_4 = \SI{100}{\ohm}$. It is
balanced at \SI{0}{\celsius} and delivers 0 to \SI{37.18}{\milli\volt} over
the measurement range. The instrumentation amplifier amplifies this to
\SI{0}{\volt} at the bottom and exactly \SI{4.00}{\volt} at the top of the
range. The following low-pass filter (chapter~6), the sample-and-hold
(chapter~7) and the 3-bit flash ADC (chapter~8) expect exactly this range:
with 8 codes over \SI{4}{\volt} one LSB is \SI{0.5}{\volt}.\newpage}

\begin{description}
\item[{\rot\bf The signal chain of the capstone}]~\\[-\lblskip]
\begin{center}
\resizebox{\linewidth}{!}{%
\begin{tikzpicture}[font=\sffamily\small, >=Stealth,
  blk/.style={draw, thick, rounded corners=1pt, minimum height=11mm, minimum width=19mm, align=center}]
  \node[blk, fill=rwucyan!10] (b) {Pt100\\bridge};
  \node[blk, fill=rwuviolet!15, right=9mm of b] (i) {2 buffers +\\diff.\ amp\\$G_1=10$};
  \node[blk, fill=rwuviolet!15, right=9mm of i] (g) {gain stage\\$G_2=10.76$\\trim};
  \node[blk, fill=rwucyan!10, right=9mm of g] (f) {low-pass\\\SI{2.34}{\hertz}\\(ch.~6)};
  \node[blk, fill=rwucyan!10, right=9mm of f] (s) {S\&H\\(ch.~7)};
  \node[blk, fill=rwucyan!10, right=9mm of s] (a) {3-bit\\flash ADC\\(ch.~8)};
  \draw[->, thick] (b) -- node[above, font=\scriptsize]{0..37\,mV} (i);
  \draw[->, thick] (i) -- (g);
  \draw[->, thick] (g) -- node[above, font=\scriptsize]{0..4\,V} (f);
  \draw[->, thick] (f) -- (s);
  \draw[->, thick] (s) -- (a);
  \draw[->, thick] (a) -- ++(1.3,0) node[right]{code};
  \node[below=2mm of b, font=\scriptsize, align=center] {50\,Hz, 5\,mV\\in series with Pt100};
\end{tikzpicture}}
\end{center}
\bi
	\ite bridge from \SI{4.00}{\volt}: $V_{BR} = \Vref\left(\frac{R}{R_1+R}-\frac{R_4}{R_3+R_4}\right) = \SIrange{0}{37.18}{\milli\volt}$ for $R = \SIrange{100}{138.5}{\ohm}$
	\ite total gain $G = \SI{4.00}{\volt}/\SI{37.18}{\milli\volt} = 107.6$
	\ite split: difference amplifier $G_1 = \SI{100}{\kilo\ohm}/\SI{10}{\kilo\ohm} = 10$, non-inverting stage $G_2 = 10.76$ -- both $<100$
	\ite zero: set by the bridge ($R_4 = R(\SI{0}{\celsius})$); span: one resistor $R_F$ in the second stage
\ei
\nsl{The common-mode voltage of this bridge is only about \SI{0.1}{\volt}
($R_4$ is small compared with $R_3$), so the op-amps must accept input
voltages close to their negative supply. With a symmetric supply or an
op-amp whose input range includes ground this is no problem.}
\os{\newpage}

\item[{\rot\bf Gain stage with trimmer}]~\\[-\lblskip]
\begin{center}
\begin{circuitikz}[opfig=1.2]
  \draw (0,0) node[op amp, noinv input up](OA){};
  \coordinate (X) at ($(OA.out)+(1.2,0)$);
  \draw (OA.+) to[short,-o] ++(-1,0) node[left]{$V_\mathrm{diff}$};
  \draw (OA.-) -- ++(0,-0.9) coordinate(N);
  \draw (N) to[R, l=\SI{91}{\kilo\ohm}] ++(1.6,0) coordinate(M) to[vR, l=$R_\mathrm{trim}$] (M -| X) -- (X);
  \draw (N) to[R, l_=\SI{10}{\kilo\ohm}] ++(0,-1.4) node[ground]{};
  \draw (OA.out) -- (X) to[short,-o] ++(0.8,0) node[right]{$\Vout$};
\end{circuitikz}
\end{center}
\bi
	\ite required: $G_2 = \SI{4.00}{\volt}/(10\cdot\SI{37.18}{\milli\volt}) = 10.758$ $\Rightarrow$ $R_F = \SI{97.58}{\kilo\ohm}$
	\ite fixed E96 value: $R_F = \SI{97.6}{\kilo\ohm}$, $G_2 = 10.76$, full scale \SI{4.0005}{\volt} (used in Falstad)
	\ite in hardware: $R_F = \SI{91}{\kilo\ohm}$ + \SI{10}{\kilo\ohm} trimmer: $G_2 = 10.1 \ldots 11.1$, full scale \SIrange{3.76}{4.13}{\volt}; setting $R_\mathrm{trim} = \SI{6.6}{\kilo\ohm}$
	\ite the fixed resistor carries most of the value: fine trimmer resolution, enough range for the tolerances
\ei
\os{\newpage}

\item[{\rot\bf Zero adjustment at the reference input}]~\\[-\lblskip]
\begin{center}
\begin{circuitikz}[opfig=1.1]
  \draw (0,0) node[op amp](D){};
  \coordinate (X) at ($(D.out)+(0.5,0)$);
  \draw (D.-) -- ++(-0.3,0) coordinate(n);
  \draw (D.+) -- ++(-0.3,0) coordinate(p);
  \draw (n) to[R, l_=$R_1$, -o] ++(-2,0) node[left]{$V_2$};
  \draw (p) to[R, l=$R_3$, -o] ++(-2,0) node[left]{$V_1$};
  \draw (n) -- ++(0,1) coordinate(t) to[R, l=$R_2$] (t -| X) -- (X);
  \draw (p) to[R, l=$R_4$] ++(0,-1.4) coordinate(r);
  \draw (D.out) -- (X) to[short,-o] ++(0.8,0) node[right]{$V_\mathrm{diff} = \frac{R_2}{R_1}(V_1-V_2) + V_\mathrm{REF}$};
  \draw (-4.5,-2.8) node[op amp, noinv input up](B){};
  \draw (B.out) -| (r);
  \draw (B.-) -- ++(0,-0.5) -| ($(B.out)+(0.25,0)$);
  \draw (B.+) to[short,-o] ++(-0.6,0) node[left]{$V_\mathrm{REF}$};
  \draw (r) node[circ]{};
  \node[right] at (r) {\scriptsize REF};
\end{circuitikz}
\end{center}
\bi
	\ite the lower end of $R_4$ is the \textbf{reference input}: $V_\mathrm{REF}$ adds directly to the output
	\ite it must be driven from a low impedance (follower), otherwise it changes $R_4$ and the CMRR
	\ite in the capstone the zero is set by the bridge ($R_4 = R(\SI{0}{\celsius})$); a REF adjustment of a few \si{\milli\volt} is needed only if op-amp offsets or resistor tolerances shift it
\ei
\os{\newpage}

\item[{\rot\bf Trimming procedure}]~\\[-\lblskip]
\bi
	\ite two adjustments: zero (offset) first, then span (gain)
\ei
\designcard{Measurement amplifier (zero and span)}{
\begin{enumerate}
\item Calculate the bridge span $V_{BR,\mathrm{max}}$ and the total gain $G = V_\mathrm{FS}/V_{BR,\mathrm{max}}$ (here $4.00\,\mathrm{V}/37.18\,\mathrm{mV} = 107.6$).
\item Split the gain: input stage (buffers + difference amplifier) $G_1 = 10$, second stage $G_2 = G/G_1 = 10.76$, each stage $G<100$.
\item Use a fixed resistor for most of $R_F$ and a trimmer for the rest (\SI{91}{\kilo\ohm} + \SI{10}{\kilo\ohm}).
\item \textbf{Zero}: set the sensor to the bottom of the range (\SI{100}{\ohm}, or a precise resistor) and check $\Vout = 0.00\,\mathrm{V}$; correct with $R_4$ or $V_\mathrm{REF}$ if necessary.
\item \textbf{Span}: set the sensor to the top of the range (\SI{138.5}{\ohm}), adjust $R_\mathrm{trim}$ until $\Vout = 4.00\,\mathrm{V}$.
\item Check the zero again.
\end{enumerate}}
\nsl{The order matters: the span trimmer changes the gain of everything that
arrives at the second stage, including a remaining offset. If the zero is
correct first, a later change of the gain no longer moves the zero point.
In the capstone the zero is fixed by the bridge and the span by $R_F$, so
the two adjustments do not influence each other.}
\os{\newpage}

\item[{\rot\bf The complete front end in Falstad}]~\\[-\lblskip]
\bi
	\ite bridge \SI{4.00}{\volt}, $R_1 = R_3 = \SI{3.9}{\kilo\ohm}$, $R_4 = \SI{100}{\ohm}$, Pt100 with a slider from \SI{100}{\ohm} to \SI{138.5}{\ohm}
	\ite two followers, difference amplifier \SI{10}{\kilo\ohm}/\SI{100}{\kilo\ohm}, gain stage $1+(\SI{91}{\kilo\ohm}+R_\mathrm{trim})/\SI{10}{\kilo\ohm}$, $R_\mathrm{trim}=\SI{6.6}{\kilo\ohm}$
	\ite readings: $V_R = \SI{0.100}{\volt}$, $V_{BR} = \SI{37.18}{\milli\volt}$, $V_\mathrm{diff}=\SI{372}{\milli\volt}$, $\Vout = \SI{4.00}{\volt}$ at \SI{138.5}{\ohm}; $\Vout=\SI{0.00}{\volt}$ at \SI{100}{\ohm}
	\ite the 50\,Hz interference of \SI{5}{\milli\volt} in series with the Pt100 is a \emph{differential} signal: \SI{4.9}{\milli\volt} at the tap, $\times 107.6$ = \SI{0.52}{\volt} amplitude
	\itee no CMRR helps here -- the low-pass filter of chapter~6 has to remove it
	\itee the output must have headroom: $\SI{4}{\volt} + \SI{0.52}{\volt}$ peak
\ei
\falstad{opamp_ina}\par
\os{\newpage}
\end{description}

%%%%%%%%%%%%%%%%%%%%%%%%%%%%%%%%%%%%%%%%%%%%%%%%%%%%%%%%%%%%%%%%%%%%%%%%%%%
\section{Exercises}
%%%%%%%%%%%%%%%%%%%%%%%%%%%%%%%%%%%%%%%%%%%%%%%%%%%%%%%%%%%%%%%%%%%%%%%%%%%

\exercise{Non-inverting and inverting amplifier}\label{ex:opamp-basic}
A sine voltage with \SI{0.5}{\volt} amplitude and \SI{1}{\kilo\hertz} is to be
amplified. Resistors: $R_1 = \SI{10}{\kilo\ohm}$, $R_2 = \SI{20}{\kilo\ohm}$.
\begin{talist}
\ta Calculate gain and output amplitude of the non-inverting amplifier.
\ta Calculate gain and output amplitude of the inverting amplifier. What is the phase shift?
\ta What input resistance does the source see in both circuits?
\ta Which resistor would you change to get an inverting gain of $-4.7$?
\ta Check a) and b) in Falstad.
\end{talist}

\begin{solution}
\begin{talist}
\ta $G = 1 + R_2/R_1 = 1 + 20/10 = 3$; output amplitude $3\cdot\SI{0.5}{\volt} = \SI{1.5}{\volt}$, in phase with the input.
\ta $G = -R_2/R_1 = -2$; output amplitude \SI{1.0}{\volt}, phase shift \ang{180}.
\ta Non-inverting: practically infinite (the op-amp input, \si{\mega\ohm} to \si{\tera\ohm}).
    Inverting: $R_1 = \SI{10}{\kilo\ohm}$, because the inverting input is a virtual ground.
\ta $R_2 = 4.7\cdot R_1 = \SI{47}{\kilo\ohm}$ (E12 value).
\ta \falstad{opamp_noninv_inv}\newline the scope shows \texttt{Vout\_ni} with \SI{\pm 1.50}{\volt} and
    \texttt{Vout\_inv} with \SI{\pm 1.00}{\volt}, inverted.
\end{talist}
\end{solution}

\exercise{Summing amplifier}\label{ex:opamp-summing}
An inverting summing amplifier has $R_f = \SI{10}{\kilo\ohm}$ and three inputs:
$V_1 = \SI{1}{\volt}$ via $R_1=\SI{10}{\kilo\ohm}$, $V_2 = \SI{2}{\volt}$ via
$R_2 = \SI{20}{\kilo\ohm}$, $V_3 = \SI{-0.5}{\volt}$ via $R_3 = \SI{10}{\kilo\ohm}$.
\begin{talist}
\ta Calculate $\Vout$ and the voltage at the inverting input.
\ta Calculate the current through $R_f$.
\ta How must $R_3$ be changed so that $V_3$ is weighted with $-0.25$?
\ta Check a) in Falstad.
\end{talist}

\begin{solution}
\begin{talist}
\ta $\Vout = -\SI{10}{\kilo\ohm}\left(\frac{\SI{1}{\volt}}{\SI{10}{\kilo\ohm}} + \frac{\SI{2}{\volt}}{\SI{20}{\kilo\ohm}} + \frac{\SI{-0.5}{\volt}}{\SI{10}{\kilo\ohm}}\right) = -(1 + 1 - 0.5)\,\si{\volt} = \SI{-1.50}{\volt}$.
    The inverting input is a virtual ground: $V_- = \SI{0}{\volt}$.
\ta $I = \SI{0.1}{\milli\ampere} + \SI{0.1}{\milli\ampere} - \SI{0.05}{\milli\ampere} = \SI{0.15}{\milli\ampere}$,
    which gives $\Vout = -\SI{10}{\kilo\ohm}\cdot \SI{0.15}{\milli\ampere} = \SI{-1.5}{\volt}$.
\ta Weight $-R_f/R_3 = -0.25$ $\Rightarrow$ $R_3 = \SI{40}{\kilo\ohm}$ (e.g.\ \SI{39}{\kilo\ohm} + \SI{1}{\kilo\ohm}).
\ta \falstad{opamp_summing}\newline \texttt{Vout} $=\SI{-1.50}{\volt}$, \texttt{V\_sum\_node} $= \SI{0.00}{\volt}$.
\end{talist}
\end{solution}

\exercise{Difference amplifier and resistor matching}\label{ex:opamp-diff}
A difference amplifier with $R_1 = R_3 = \SI{10}{\kilo\ohm}$ and $R_2 = R_4 = \SI{100}{\kilo\ohm}$
amplifies a bridge signal: $V_1 = \SI{2.505}{\volt}$, $V_2 = \SI{2.500}{\volt}$.
\begin{talist}
\ta Calculate differential and common-mode voltage and the output voltage.
\ta $R_4$ is \SI{1}{\percent} too high (\SI{101}{\kilo\ohm}). Calculate the output voltage.
\ta Calculate the common-mode gain and the CMRR for b).
\ta Estimate the worst-case CMRR with \SI{1}{\percent} resistors.
\ta Check a) and b) in Falstad.
\end{talist}

\begin{solution}
\begin{talist}
\ta $V_D = \SI{5}{\milli\volt}$, $V_{CM} = \SI{2.5025}{\volt}$; $\Vout = 10\cdot\SI{5}{\milli\volt} = \SI{50.0}{\milli\volt}$.
\ta $V_+ = \SI{2.505}{\volt}\cdot\frac{101}{111} = \SI{2.27932}{\volt}$,
    $\Vout = V_+\,(1+10) - \SI{2.5}{\volt}\cdot 10 = \SI{25.0726}{\volt}-\SI{25}{\volt} = \SI{72.6}{\milli\volt}$:
    \SI{45}{\percent} too high.
\ta $A_{CM} = 11\cdot\frac{101}{111} - 10 = 0.0090$ (input both at \SI{2.5}{\volt} gives \SI{22.5}{\milli\volt});
    $A_D\approx 10$; $\mathrm{CMRR} = 10/0.009 \approx 1100$ (\SI{61}{\deci\bel}).
\ta $\mathrm{CMRR} \approx (1+G)/(4\varepsilon) = 11/0.04 = 275$ (\SI{49}{\deci\bel}): with
    \SI{2.5}{\volt} common mode up to \SI{9}{\milli\volt} error at the input, i.e.\ almost half of the
    full bridge span. Matched resistors or an integrated difference amplifier are needed.
\ta \falstad{opamp_diffamp}\newline \texttt{Vout\_matched} $= \SI{50.0}{\milli\volt}$, \texttt{Vout\_mismatch} $= \SI{72.6}{\milli\volt}$.
\end{talist}
\end{solution}

\exercise{Finite open-loop gain}\label{ex:opamp-gain}
A signal of \SI{1}{\milli\volt} DC is amplified with non-inverting stages
($R_1 = \SI{1}{\kilo\ohm}$). The op-amps have $A = 10^5$ (as in Falstad).
\begin{talist}
\ta Calculate $\Vout$ for one stage with $G=1000$ and one stage with $G=10000$. Give $R_2$ and the errors.
\ta Calculate $\Vout$ for $G = 10\cdot 100$ and $G = 100\cdot 100$ in two stages.
\ta How large is the error of the $G=100$ stage at \SI{1}{\kilo\hertz} if the op-amp has $\mathrm{GBW} = \SI{1}{\mega\hertz}$?
\ta Check a) and b) in Falstad.
\end{talist}

\begin{solution}
\begin{talist}
\ta $G=1000$: $R_2 = \SI{999}{\kilo\ohm}$, $G_\mathrm{real} = 1000/(1+0.01) = 990.1$,
    $\Vout = \SI{0.990}{\volt}$ (\SI{-1}{\percent}).
    $G=10000$: $R_2 = \SI{9.999}{\mega\ohm}$, $G_\mathrm{real} = 10000/1.1 = 9091$, $\Vout = \SI{9.09}{\volt}$ (\SI{-9.1}{\percent}).
\ta $10/(1+10^{-4})\cdot 100/(1+10^{-3}) = 9.999\cdot 99.90 = 998.9$: $\Vout = \SI{0.999}{\volt}$ (\SI{-0.11}{\percent}).
    $(100/1.001)^2 = 9980$: $\Vout = \SI{9.98}{\volt}$ (\SI{-0.2}{\percent}).
\ta $A(\SI{1}{\kilo\hertz}) = \SI{1}{\mega\hertz}/\SI{1}{\kilo\hertz} = 1000$, error $\approx G/A = \SI{10}{\percent}$
    (and a phase shift). Falstad does not show this, because it does not simulate the GBW.
\ta \falstad{opamp_finite_gain}\newline \texttt{Vout\_1000} $=\SI{0.990}{\volt}$, \texttt{Vout\_10x100} $=\SI{0.999}{\volt}$,\newline
    \texttt{Vout\_10000} $=\SI{9.09}{\volt}$, \texttt{Vout\_100x100} $=\SI{9.98}{\volt}$.
\end{talist}
\end{solution}

\exercise{Offset voltage and bias current}\label{ex:opamp-offset}
A non-inverting amplifier has $R_1 = \SI{1}{\kilo\ohm}$, $R_2 = \SI{100}{\kilo\ohm}$; its input is at \SI{0}{\volt}.
\begin{talist}
\ta The op-amp has an offset voltage of \SI{1}{\milli\volt}. Calculate $\Vout$.
\ta The input is connected to ground through $R_S = \SI{100}{\kilo\ohm}$ and a bias current
    of \SI{100}{\nano\ampere} flows into the non-inverting input (no offset). Calculate $\Vout$.
\ta The capstone amplifier has $G = 107.6$ and the op-amps have \SI{2}{\milli\volt} offset (referred to the input).
    Compare the output error with the full scale of \SI{4}{\volt} and with one LSB of the 3-bit ADC.
\ta Check a) and b) in Falstad.
\end{talist}

\begin{solution}
\begin{talist}
\ta $\Vout = (1 + 100)\cdot\SI{1}{\milli\volt} = \SI{0.101}{\volt}$.
\ta $V_+ = -R_S I_B = -\SI{100}{\kilo\ohm}\cdot\SI{100}{\nano\ampere} = \SI{-10}{\milli\volt}$, $\Vout = 101\cdot(\SI{-10}{\milli\volt}) = \SI{-1.01}{\volt}$.
    A CMOS op-amp with $I_B=\SI{1}{\pico\ampere}$ would give only \SI{10}{\micro\volt}.
\ta $107.6\cdot\SI{2}{\milli\volt} = \SI{0.215}{\volt}$: \SI{5.4}{\percent} of the full scale, 0.43 LSB
    ($\LSB = \SI{4}{\volt}/8 = \SI{0.5}{\volt}$). The zero must be checked and trimmed, or a zero-drift op-amp used.
\ta \falstad{opamp_offset}\newline \texttt{Vout\_offset} $=\SI{0.101}{\volt}$, \texttt{V\_plus\_b} $=\SI{-10.0}{\milli\volt}$, \texttt{Vout\_bias} $=\SI{-1.01}{\volt}$.
\end{talist}
\end{solution}

\exercise{Capstone: instrumentation amplifier with 4.00\,V full scale}\label{ex:opamp-ina}
The Pt100 of the capstone (\SI{100}{\ohm} at \SI{0}{\celsius}, \SI{138.5}{\ohm} at \SI{100}{\celsius})
sits in a bridge fed from the \SI{4.00}{\volt} reference: $R_1$ (\SI{3.9}{\kilo\ohm}) above the
Pt100 gives $V_S$, $R_3 = \SI{3.9}{\kilo\ohm}$ above $R_4 = \SI{100}{\ohm}$ gives $V_R$.
The amplifier consists of two followers, a difference amplifier and a non-inverting stage.
Its output shall be \SI{0}{\volt} at \SI{0}{\celsius} and exactly \SI{4.00}{\volt} at \SI{100}{\celsius}.
\begin{talist}
\ta Calculate $V_R$, $V_S$ and the bridge output $V_{BR} = V_S - V_R$ at \SI{0}{\celsius} and \SI{100}{\celsius}. What is the common-mode voltage?
\ta Calculate the total gain. Why is it split into two stages?
\ta Choose $G_1 = 10$ for the difference amplifier and give its resistors.
\ta Dimension the second stage with $R_g = \SI{10}{\kilo\ohm}$. Give the exact $R_F$, the nearest E96 value and a version with \SI{91}{\kilo\ohm} plus a \SI{10}{\kilo\ohm} trimmer (setting and full-scale range).
\ta The three op-amp INA could replace the followers. Which $R_G$ gives an input stage gain of 10 with $R = \SI{10}{\kilo\ohm}$?
\ta A 50\,Hz interference of \SI{5}{\milli\volt} amplitude is in series with the Pt100. How large is it at the output? What follows for the supply voltage and the next stage?
\ta Check the design in Falstad.
\end{talist}

\begin{solution}
\begin{talist}
\ta $V_R = \SI{4}{\volt}\cdot 100/4000 = \SI{0.100}{\volt}$. \SI{0}{\celsius}: $V_S = \SI{0.100}{\volt}$, $V_{BR} = 0$
    (sensor current \SI{1.000}{\milli\ampere}).
    \SI{100}{\celsius}: $V_S = \SI{4}{\volt}\cdot 138.5/4038.5 = \SI{0.13718}{\volt}$, $V_{BR} = \SI{37.18}{\milli\volt}$.
    The common mode is only about \SI{0.1}{\volt}: the op-amp inputs must work close to the negative rail.
\ta $G = \SI{4.00}{\volt}/\SI{37.18}{\milli\volt} = 107.6$. One stage would be possible (error $G/A = \SI{0.11}{\percent}$),
    but two stages of about 10 keep each error at \SI{0.01}{\percent} and the bandwidth high, follow the lab rule $G<100$,
    and separate zero (bridge) and span ($R_F$).
\ta $R_1 = R_3 = \SI{10}{\kilo\ohm}$, $R_2 = R_4 = \SI{100}{\kilo\ohm}$ (\SI{0.1}{\percent} or matched).
    $V_\mathrm{diff} = 10\cdot\SI{37.18}{\milli\volt} = \SI{371.8}{\milli\volt}$ at \SI{100}{\celsius}.
\os{\newpage}
\ta $G_2 = \SI{4.00}{\volt}/\SI{371.8}{\milli\volt} = 10.758$, $R_F = 9.758\cdot\SI{10}{\kilo\ohm} = \SI{97.58}{\kilo\ohm}$.
    E96: \SI{97.6}{\kilo\ohm}, $G_2 = 10.76$, full scale \SI{4.0005}{\volt}.
    Trimmer version: \SI{91}{\kilo\ohm} + $R_\mathrm{trim} = \SI{6.6}{\kilo\ohm}$; range $G_2 = 10.1\ldots 11.1$,
    full scale \SIrange{3.76}{4.13}{\volt}.
\ta $1+2R/R_G = 10$ $\Rightarrow$ $R_G = 2R/9 = \SI{2.22}{\kilo\ohm}$; then the difference stage can have gain 1 and the
    second stage stays as it is.
\ta At the tap: $\SI{5}{\milli\volt}\cdot\frac{\SI{3.9}{\kilo\ohm}}{\SI{3.9}{\kilo\ohm}+\SI{100}{\ohm}} = \SI{4.875}{\milli\volt}$,
    at the output $\times 107.6 = \SI{0.52}{\volt}$ amplitude: the output swings up to \SI{4.52}{\volt}.
    The op-amps need headroom (e.g.\ \SI{\pm 15}{\volt} as in Falstad); the low-pass of
    chapter~6 must attenuate the interference below $\LSB/2 = \SI{0.25}{\volt}$.
\ta \falstad{opamp_ina}\newline with the Pt100 at \SI{138.5}{\ohm} the labels show \texttt{VB\_minus} $= \SI{0.100}{\volt}$,
    \texttt{V\_diff} $= \SI{372}{\milli\volt}$ and \texttt{Vout} $= \SI{4.00}{\volt}$; with the slider at \SI{100}{\ohm} \texttt{Vout} $= \SI{0.00}{\volt}$.
\end{talist}
\end{solution}

%%% Local Variables:
%%% mode: latex
%%% TeX-master: "docu"
%%% End:
